% !TEX root = ../../main.tex
\subsection{系统总体架构}

FoodShield 采用多端协同的分层架构，整体上可以划分为表现层、业务服务层、安全机制层、数据存储层和审计验证层。表现层包括用户端、骑手端和管理员端；业务服务层负责注册登录、创建订单、订单历史、订单搜索、备注标签、订单认证、消息转发、投诉申请等功能；安全机制层提供 PID 生成、Token 验证、SM4 加密、SM3 摘要和 Merkle Tree 构建等能力；数据存储层保存用户、骑手、订单、消息、快照和审计日志等数据；审计验证层负责日志完整性验证、Root 快照比对、条件溯源记录和三端摘要一致性增强。

系统总体架构可以概括如下：

\begin{enumerate}[label=\arabic*., leftmargin=2em]
\item \textbf{用户端}：用户端提供用户注册、登录、创建订单、查看订单历史、搜索订单、添加备注标签、进入订单通信和提交投诉申请等功能。用户完成注册后，系统在内部生成 PID，并在订单创建时将 PID 与订单号进行绑定。用户进入通信通道前，需要提交订单号和 Token 完成订单认证。

\item \textbf{骑手端}：骑手端提供骑手登录、查看分配订单、根据订单号搜索订单、进入订单通信和提交异常申请等功能。骑手端只显示订单相关必要信息，不显示用户真实身份和 PID，避免骑手通过长期标识进行跨订单关联。

\item \textbf{平台服务器}：平台服务器是系统核心部分，负责接收用户端、骑手端和管理员端请求，并调用不同安全模块完成业务处理。平台服务器内部包括用户注册与机器识别模块、匿名身份管理模块、订单认证模块、消息转发模块、消息加密存储模块、日志审计模块和条件溯源模块。

\item \textbf{管理员端}：管理员端用于平台审计和异常处理。管理员登录后可以根据订单号检索订单摘要，查看消息哈希和 Merkle Root，执行日志完整性验证，处理用户端或骑手端提交的溯源申请。管理员端默认不展示通信明文内容，也不展示用户备注或标签原文。

\item \textbf{审计日志系统}：审计日志系统保存通信消息摘要、Merkle Root 快照和管理员操作日志。系统在订单通信过程中持续生成消息哈希，并在订单结束或管理员触发快照时计算 Merkle Root。后续验证时，系统重新计算 Root 并与快照值对比，以判断日志是否被篡改。

\item \textbf{三端摘要一致性增强机制}：在订单通信结束时，系统可将订单通信摘要或 Merkle Root 摘要分别留存在用户端、骑手端和平台端。后续审计时，若攻击者试图篡改平台侧数据库记录，则需要同时修改三端摘要才能伪造一致结果，从而提高日志篡改成本。
\end{enumerate}

系统的基本业务流程如下。首先，用户完成注册和登录，注册阶段包含用户名、口令和机器识别信息，平台服务器生成用户 PID，并保存 PID 与真实身份之间的受控映射关系。其次，用户创建订单，服务器生成订单号并将订单与用户 PID、骑手身份、订单备注和订单状态进行绑定，同时生成订单 Token。随后，用户进入订单通信页面时提交订单号和 Token，服务器验证通过后允许用户加入对应 WebSocket 通信房间；骑手端经过身份验证后进入同一订单房间。通信过程中，平台服务器负责消息转发、消息加密存储和消息哈希计算。订单完成或管理员生成审计快照时，系统根据消息摘要构建 Merkle Tree 并保存 Merkle Root。若后续发生投诉或纠纷，用户或骑手可基于订单号发起申请，管理员经权限验证后查看审计摘要、验证日志完整性，并在必要时执行条件溯源。

在架构设计上，FoodShield 的关键特点在于将业务流程和安全机制进行绑定。订单不是单纯的业务编号，而是认证、通信、日志和溯源的核心索引；PID 不是公开展示的用户身份，而是内部匿名绑定标识；Token 不是普通随机字符串，而是基于 HMAC-SM3 的订单认证凭证；消息记录不是普通数据库文本，而是加密存储并参与 Merkle Tree 完整性验证的审计对象；管理员端不是明文查看入口，而是哈希摘要、Root 验证和条件溯源的受控审计入口。通过这种设计，系统能够围绕订单这一核心业务对象建立完整的隐私保护与可信审计闭环。
